\documentclass[11pt,a4paper]{article}
\usepackage[utf8]{inputenc}\usepackage[T1]{fontenc}
\usepackage{amsmath,amssymb,bm,array,geometry,color,fancyhdr,hyperref}
\geometry{left=19mm,right=19mm,top=18mm,bottom=19mm}
\definecolor{ink}{RGB}{20,52,72}\definecolor{accent}{RGB}{22,115,139}
\hypersetup{colorlinks=true,urlcolor=accent,linkcolor=ink}
\setlength{\parindent}{0pt}\setlength{\parskip}{3pt}
\setlength{\abovedisplayskip}{6pt}\setlength{\belowdisplayskip}{6pt}\setlength{\jot}{3pt}
\pagestyle{fancy}\fancyhf{}\renewcommand{\headrulewidth}{0pt}
\fancyfoot[L]{\footnotesize pysnspd 1.1.3 / candidato / \href{Supuestos.pdf}{Supuestos}}\fancyfoot[R]{\small\thepage}
\newcommand{\block}[1]{\par\vspace{4pt}{\large\bfseries\color{ink}#1}\par\vspace{3pt}}
\newcommand{\labelrow}[2]{\(\displaystyle #1\)& #2\\}
\newcommand{\origin}[1]{\par{\footnotesize\textbf{Campo / procedencia:} #1}\par}
\newcommand{\code}[1]{\par{\footnotesize\textbf{Código:} \texttt{#1}}\par}
\newcommand{\dd}{\,\mathrm d}\newcommand{\FD}{f_{\rm FD}}\newcommand{\BE}{n_{\rm BE}}
\newcommand{\Rea}{\operatorname{Re}}\newcommand{\Ima}{\operatorname{Im}}
\newcommand{\refs}{\vfill{\footnotesize\href{Supuestos.pdf}{Hipótesis y límites: Supuestos}; \href{Parametros.pdf}{entradas: Parametros}.}}
\begin{document}
{\LARGE\bfseries\color{ink}Ecuaciones principales en 2D}\par
{\large Condensado, cuasipartículas y fonones: candidato 1.1.3}\par
{\small BCS/KWT efectivo con corrección térmica de depareo anclada a DC. Candidato: validación dinámica pendiente.}
\renewcommand{\theequation}{1.\arabic{equation}}
\hypertarget{G-final}{}\block{Condensado: evolución KWT}
\begin{gather}
E_u=k_BT_c,\quad d_i=\Delta_i/E_u,\quad \rho_i=|d_i|^2,\quad
 E_{ik}=E_u e_{ik},\quad e_{ik}=\sqrt{\bar\xi_k^2+\rho_i},\\
 D_td_i=\dot d_i+\frac{2ie}{\hbar}\phi_i d_i,\qquad
 z_i=\begin{cases}d_i/|d_i|,&d_i\ne0,\\1,&d_i=0,\end{cases}\\
\mathcal F_i=
 \partial_{d_i^*}U-\sum_jT_{E,j}\partial_{d_i^*}S_j
 =\left.\partial_{d_i^*}(U_{\rm BCS}+G)\right|_{p,\mathbf T_E,\alpha},
 \quad f_i=z_i^*\mathcal F_i/u_i^*,\\
\boxed{D_td_i=-\frac{z_i}{A_i\tau_K}
 \left(\frac{\Rea f_i}{R_i}+iR_i\Ima f_i\right)},\\
u_i^*=N_0E_u^2V_i,\quad
 A_i=\sqrt{\frac{1+\vartheta_i}{2}},\quad
 \vartheta_i=\frac{\max(T_{E,i},T_b)}{T_c},\quad
 \tau_K=\frac{\pi\hbar}{8k_BT_c},\\
\tau_{{\rm sc},i}^{-1}
 =\frac{\vartheta_i}{\tau_{ee}(T_c)}
 +\frac{\vartheta_i^3}{\tau_{ep}(T_c)},\quad
 R_i=\sqrt{1+(2E_u\tau_{{\rm sc},i}|d_i|/\hbar)^2},\\
\boxed{Q_{\Delta,i}
 =\frac{2u_i^*}{A_i\tau_K}
 \left[\frac{(\Rea f_i)^2}{R_i}+R_i(\Ima f_i)^2\right]\geq0}.
\end{gather}
\par\begin{tabular}{@{}p{.29\linewidth}p{.67\linewidth}@{}}
\hline
\labelrow{i,k,t,\mathbf r_i}{nodo espacial, estado nativo, tiempo [s], posición [m]}
\labelrow{\Delta_i,d_i,\rho_i,z_i}{gap [J], gap normalizado, módulo al cuadrado, fase}
\labelrow{E_{ik},\bar\xi_k}{excitación [J], energía normal nativa dividida por \(E_u\)}
\labelrow{U,S_j,G,\mathcal F_i}{energía [J], entropía nodal [J/K], energía libre de gradiente [J], fuerza [J]}
\labelrow{\phi_i,e,\hbar,k_B}{potencial [V], carga positiva [C], acción [J s], constante [J/K]}
\labelrow{N_0,V_i,u_i^*}{DOS monoespín [J\(^{-1}\)m\(^{-3}\)], volumen [m\(^{3}\)], escala [J]}
\labelrow{T_E,T_b,T_c}{temperaturas equivalente electrónica, baño y crítica [K]}
\labelrow{A_i,R_i,f_i,Q_{\Delta,i}}{factores y fuerza adimensionales; calor KWT [W]}
\labelrow{\tau_{ee}(T_c),\tau_{ep}(T_c)}{calibradores libres Allmaras [s], no constantes universales}
\hline
\end{tabular}\par
{\small En la primera igualdad de $\mathcal F_i$, las derivadas de $U,S_j$ siguen $T_E(d,p)$ a $p,\alpha$ fijos, con $S_j=S_{qp,j}-G_{T_j}$. En la última se mantiene además $\mathbf T_E$ fijo; las contribuciones térmicas se cancelan por la regla de la cadena. Definición completa de $G$: \href{Auxiliar.pdf\#nameddest=G-final}{Auxiliar, p.~1}.}
\origin{BCS + disipación KWT. \href{https://doi.org/10.1103/PhysRevLett.40.1041}{Kramer--Watts-Tobin (1978)}; \href{https://thesis.caltech.edu/13748/}{Allmaras (2020)}, ec.\ 3.24. Matching propio: \(H_i^{\mathrm{match}}\ge0\) persiste al elevar \(T_E\), es importante en el hotspot y limita la predicción de latencia; \href{Supuestos.pdf\#nameddest=S15}{S15}, \href{Auxiliar.pdf}{Auxiliar}.}
\refs
\newpage
\block{Poblaciones electrónicas y fonónicas}
\begin{gather}
\boxed{\begin{aligned}\dot p_{ik}={}&
 \frac{1}{4w_k}\sum_a\nu^e_{ka}r_{ia}
 -\frac{1}{4w_kV_i}\sum_{\ell,a}B_{i\ell}\beta^{(i)}_{ka}J_{\ell a}\\
 &+\frac{\FD(E_{ik},T_{E,i})-p_{ik}}{\tau_{{\rm BGK},i}}
 +h_{ik}P_i
 \end{aligned}},\\
\boxed{\dot n_{iq}=
 \frac1{c_q}\sum_a\nu^{ph}_{qa}r_{ia}
 -\frac{n_{iq}-\BE(\Omega_q,T_b)}{\tau_{{\rm esc},i}}},\\
\tau_{{\rm BGK},i}=\tau_{ee}(T_c)\frac{T_c}{\max(T_{E,i},T_b)},
 \quad c_q=\frac{g_{ph}(\Omega_q)W_q^{ph}}{N_0E_u},
 \quad \Omega_q=E_u\omega_q,\\
\boxed{P_i=Q_{\Delta,i}+P_{J,i}+T_{E,i}\dot G_{T_i}},\quad
 G_{T_i}\equiv\partial_{T_{E,i}}G=\frac{u_i^*}{T_c}\varphi_\theta(\rho_i,\theta_i,\mathcal Q_i),\\
 \epsilon_{ik}=4u_i^*w_ke_{ik},\quad
 \sigma_{ik}=\frac{\epsilon_{ik}k_B}{E_{ik}}
       \log\frac{1-p_{ik}}{p_{ik}},\\
 h_{ik}\geq0,\quad\sum_k\epsilon_{ik}h_{ik}=1,\quad
 P_i\!\left(\sum_k\sigma_{ik}h_{ik}-T_{E,i}^{-1}\right)\geq0,\\
 \sum_kw_ke_{ik}\FD(E_{ik},T_{E,i})
 =\sum_kw_ke_{ik}p_{ik},\qquad
 0\leq p_{ik}\leq1,\quad n_{iq}\geq0.
\end{gather}
{\small \(p,n\) evolucionan fuera de equilibrio. \(T_E\) es su momento energético, no una tercera población térmica. El término reversible \(T_E\dot G_T\) puede tener ambos signos. El portador usa actividades Pauli dependientes del signo de $P_i$; fuera de FD admite producción entrópica positiva. Su cálculo algebraico y el portador \(h\) figuran en \href{Auxiliar.pdf}{Auxiliar}; nunca se obtiene del residuo de energía.}
\par\begin{tabular}{@{}p{.29\linewidth}p{.67\linewidth}@{}}
\hline
\labelrow{p_{ik},n_{iq}}{ocupación de cuasipartículas por espín y de fonones}
\labelrow{k,q,a,\ell}{estado electrónico, modo fonónico, evento, cara espacial}
\labelrow{w_k,c_q,\omega_q}{pesos nativos y energía fonónica normalizada}
\labelrow{r_{ia},\nu^e,\nu^{ph}}{tasa de evento [s\(^{-1}\)]; estequiometrías compartidas}
\labelrow{B_{i\ell},\beta^{(i)}_{ka}}{incidencia orientada e interpolación a energía común}
\labelrow{J_{\ell a}}{flujo de transporte nativo [m\(^{3}\)/s]}
\labelrow{h_{ik},P_i,P_{J,i}}{portador [J\(^{-1}\)]; potencia electrónica y Joule [W]}
\labelrow{L_i,\theta_i,\mathcal Q_i,G_{T_i}}{peso [J], \(T_{E,i}/T_c\), depareo efectivo/\(E_u\), derivada [J/K]}
\labelrow{\epsilon_{ik},\sigma_{ik}}{momentos de energía [J] y entropía [J/K]}
\labelrow{\tau_{\rm esc},g_{ph},W_q^{ph}}{escape [s], DOS fonónica [J\(^{-1}\)m\(^{-3}\)], peso [J]}
\hline
\end{tabular}\par
\origin{Cinética cuántica BCS/electrón--fonón, transporte conservativo y BGK energético. \href{https://arxiv.org/pdf/1611.06060}{Vodolazov (2017)}, ecs.\ 3--4; discretización nativa y extensión termoenergética propias.}
\code{thermal\_coupled.py; resummed\_thermal.py; effective\_bcs\_kinetics.py}
\refs\end{document}
